
\documentclass[11pt,a4paper]{article}

% -------------------- Packages --------------------
\usepackage[margin=1in]{geometry}
\usepackage[T1]{fontenc}
\usepackage{lmodern}
\usepackage{microtype}
\usepackage{amsmath,amssymb}
\usepackage{booktabs}
\usepackage{longtable}
\usepackage{array}
\usepackage{xcolor}
\usepackage{listings}
\usepackage{hyperref}
\usepackage{fancyhdr}
\usepackage{enumitem}
\usepackage{float}

% -------------------- Formatting --------------------
\hypersetup{
    colorlinks=true,
    linkcolor=blue!50!black,
    urlcolor=blue!50!black,
    citecolor=blue!50!black
}

\pagestyle{fancy}
\fancyhf{}
\lhead{ScrollSense -- Assignment 2}
\rhead{Database Systems}
\cfoot{\thepage}
\setlength{\headheight}{14pt}

\definecolor{codebg}{RGB}{247,247,247}
\definecolor{codeframe}{RGB}{210,210,210}
\definecolor{codeblue}{RGB}{30,70,140}
\definecolor{codegreen}{RGB}{30,110,70}

\lstdefinestyle{sqlstyle}{
    language=SQL,
    basicstyle=\ttfamily\small,
    backgroundcolor=\color{codebg},
    frame=single,
    rulecolor=\color{codeframe},
    breaklines=true,
    showstringspaces=false,
    keywordstyle=\color{codeblue}\bfseries,
    commentstyle=\color{codegreen},
    numbers=none,
    tabsize=4,
    columns=fullflexible,
    keepspaces=true
}

\lstdefinestyle{textstyle}{
    basicstyle=\ttfamily\footnotesize,
    backgroundcolor=\color{codebg},
    frame=single,
    rulecolor=\color{codeframe},
    breaklines=true,
    numbers=none,
    columns=fullflexible,
    keepspaces=true
}

\newcommand{\studentname}{Abhijeet Kumar}
\newcommand{\studentid}{DA24B001}
\newcommand{\course}{Data Management}
\newcommand{\assignment}{Assignment 2: Query Processing and Concurrency}

\title{
    \textbf{ScrollSense}\\[0.4em]
    \large \assignment
}
\author{
    \studentname\\
    \studentid\\
    Indian Institute of Technology Madras
}
\date{October 2026}

\begin{document}

\maketitle
\thispagestyle{empty}

\vspace{1em}

\begin{abstract}
In this assignment, I investigated query execution plans, composite
indexes, join algorithms, predicate pushdown, and transaction
concurrency using the ScrollSense database. I scaled the dataset,
examined SQLite execution plans before and after index creation,
estimated the I/O cost of alternative join strategies, and rewrote
a query to make predicate pushdown explicit. I also tested SQLite's
writer-locking behaviour using two independent database connections.

The experiments show that appropriate indexes can change SQLite's
access paths, while the query planner still uses temporary structures
for operations such as grouping and ordering. The concurrency test
confirms that \texttt{BEGIN IMMEDIATE} prevents a second connection
from acquiring a writer reservation while the first connection holds
it. Where comparable timing measurements or experimental outcomes
were not recorded, I distinguish the available evidence from
theoretical analysis rather than reporting unsupported results.
\end{abstract}

\newpage
\tableofcontents
\newpage

% ============================================================
\section{Experimental Setup}
% ============================================================

\subsection{Environment}

I carried out the experiments on a MacBook Air using Python and
SQLite. I kept the original Assignment 1 files separate from the
working files for Assignment 2 so that the original schema and
queries remained available for reference.

\begin{table}[H]
\centering
\begin{tabular}{ll}
\toprule
\textbf{Item} & \textbf{Configuration} \\
\midrule
Database engine & SQLite 3.54.0 \\
Python & 3.14.8 in a virtual environment \\
Operating system & macOS \\
Database & \texttt{scrollsense\_a2.db} \\
Scale factor & 50$\times$ \\
Assumed page size for cost model & 4096 bytes \\
Assumed average row size & 64 bytes \\
\bottomrule
\end{tabular}
\caption{Experimental environment and cost-model assumptions.}
\end{table}

I generated a scaled database with a scale factor of 50. The resulting
database passed SQLite's integrity check. The following row counts
were recorded after generation.

\begin{table}[H]
\centering
\begin{tabular}{lr}
\toprule
\textbf{Table} & \textbf{Recorded rows} \\
\midrule
User & 250,000 \\
Video & 1,000,000 \\
Impression & 15,000,000 \\
ViewSegment & 8,892,547 \\
Signal & 562,275 \\
Session & 100,000 \\
Turn & 276,011 \\
ToolCall & 218,196 \\
\bottomrule
\end{tabular}
\caption{Selected row counts in the scaled database.}
\end{table}

I also made a small correction to the data generator so that generated
phone numbers remain unique at the larger scale. I kept this change
in the Assignment 2 working copy rather than modifying the original
Assignment 1 version.

\subsection{Method}

For query-plan analysis, I used SQLite's \texttt{EXPLAIN QUERY PLAN}.
I compared the recorded plans before and after creating the indexes
in \texttt{indexes.sql}. I interpreted operators such as
\texttt{SCAN}, \texttt{SEARCH}, \texttt{USING COVERING INDEX}, and
\texttt{USE TEMP B-TREE} to understand how SQLite accessed and
processed the data.

The page-transfer calculations in this report use a simplified
textbook model. They are estimates, not measurements of physical
disk I/O. Similarly, an execution plan establishes the access path
chosen by SQLite but does not, by itself, establish a runtime
improvement.

% ============================================================
\section{Task 1: Access Paths and Scan-Cost Estimate}
% ============================================================

\subsection{F4: Like Retractions Within 60 Seconds}

F4 finds cases where a user liked a video and subsequently retracted
that like within 60 seconds. The query joins \texttt{Signal} to itself.
The first alias represents the original like, and the second alias
represents the retraction.

The relevant conditions are:

\begin{lstlisting}[style=sqlstyle]
WHERE l.signal_type = 'like'
  AND r.signal_type = 'like_retraction'
  AND r.created_at > l.created_at
  AND (julianday(r.created_at) -
       julianday(l.created_at)) * 86400 <= 60
\end{lstlisting}

The query also orders the result by the original like timestamp and
signal ID.

\subsubsection{Plan before index creation}

\begin{lstlisting}[style=textstyle]
|--SCAN r
|--SEARCH l USING INTEGER PRIMARY KEY (rowid=?)
`--USE TEMP B-TREE FOR ORDER BY
\end{lstlisting}

The plan scans the retraction-side alias \texttt{r}. For each candidate,
SQLite searches for the corresponding original signal using its integer
primary key. A temporary B-tree is used to satisfy the ordering.

The full scan is important because \texttt{Signal} contains 562,275
rows in the scaled database. The plan does not mean that every pair
of signals is compared; the primary-key lookup restricts the lookup
for the matching original signal.

\subsubsection{Plan after index creation}

I created the following index:

\begin{lstlisting}[style=sqlstyle]
CREATE INDEX IF NOT EXISTS
    idx_signal_type_retracts_created
ON Signal(signal_type, retracts_signal_id, created_at);
\end{lstlisting}

The recorded plan became:

\begin{lstlisting}[style=textstyle]
|--SEARCH r USING COVERING INDEX
|  idx_signal_type_retracts_created (signal_type=?)
|--SEARCH l USING INTEGER PRIMARY KEY (rowid=?)
`--USE TEMP B-TREE FOR ORDER BY
\end{lstlisting}

SQLite now searches the index for the requested signal type rather
than scanning the entire retraction-side table. The index also contains
the referenced signal ID and timestamp, which makes it covering for
the columns needed from the \texttt{r} alias. The sort remains because
the index does not directly provide the complete requested ordering.

The F4 query took 0.09 seconds before the index was created
and 0.03 seconds after it was recreated. This is a single-run
observation rather than a controlled benchmark, so cache state
and system load may affect the measured difference.

\subsection{F7: Session-Level Inference Cost}

F7 calculates the cost of each agent session, grouped by prompt
template version. It joins \texttt{Session}, \texttt{Turn},
\texttt{PromptTemplateVersion}, and \texttt{ModelPricingPeriod}.
The query filters sessions to August 2026, aggregates token-based
costs, applies a \texttt{HAVING} condition, and sorts by session cost.

The date predicate is:

\begin{lstlisting}[style=sqlstyle]
WHERE s.started_at >= '2026-08-01T00:00:00Z'
  AND s.started_at <  '2026-09-01T00:00:00Z'
\end{lstlisting}

\subsubsection{Plan before index creation}

\begin{lstlisting}[style=textstyle]
|--SCAN t
|--SEARCH s USING INTEGER PRIMARY KEY (rowid=?)
|--SEARCH ptv USING INTEGER PRIMARY KEY (rowid=?)
|--SEARCH p USING INDEX
|  sqlite_autoindex_ModelPricingPeriod_1
|  (model_id=? AND valid_from<?)
|--USE TEMP B-TREE FOR GROUP BY
`--USE TEMP B-TREE FOR ORDER BY
\end{lstlisting}

The original plan starts by scanning \texttt{Turn}. It then looks up
the associated session, template version, and pricing period. SQLite
uses temporary B-trees for grouping and ordering.

\subsubsection{Plan after index creation}

I created an index on the session start time and session ID:

\begin{lstlisting}[style=sqlstyle]
CREATE INDEX IF NOT EXISTS
    idx_session_started_session
ON Session(started_at, session_id);
\end{lstlisting}

The recorded plan became:

\begin{lstlisting}[style=textstyle]
|--SEARCH s USING COVERING INDEX
|  idx_session_started_session
|  (started_at>? AND started_at<?)
|--SEARCH t USING INDEX sqlite_autoindex_Turn_1
|  (session_id=?)
|--SEARCH ptv USING INTEGER PRIMARY KEY (rowid=?)
|--SEARCH p USING INDEX
|  sqlite_autoindex_ModelPricingPeriod_1
|  (model_id=? AND valid_from<?)
|--USE TEMP B-TREE FOR GROUP BY
`--USE TEMP B-TREE FOR ORDER BY
\end{lstlisting}

The planner now begins with the date-filtered sessions and then
looks up matching turns. The leading \texttt{started\_at} column
supports the range predicate. The remaining joins and temporary
grouping and sorting structures are still required.

The full F7 query took 0.13 seconds before the target index
was created and 0.08 seconds after it was recreated. These are
single-run wall-clock measurements; cache state and system load
may affect the observed difference.

\subsection{F10: Rolling Watch Analysis}

F10 joins the high-volume \texttt{ViewSegment} and \texttt{Impression}
tables to analyse rolling watch behaviour. The captured plan was:

\begin{lstlisting}[style=textstyle]
|--CO-ROUTINE rolling_watch
|  |--MATERIALIZE daily_watch
|  |  |--SCAN vs
|  |  |--SEARCH i USING INTEGER PRIMARY KEY (rowid=?)
|  |  `--USE TEMP B-TREE FOR GROUP BY
|  |--SCAN d1
|  |--BLOOM FILTER ON d2 (user_id=?)
|  |--SEARCH d2 USING AUTOMATIC COVERING INDEX
|  |  (user_id=?)
|  `--USE TEMP B-TREE FOR GROUP BY
|--SCAN rolling_watch
`--USE TEMP B-TREE FOR ORDER BY
\end{lstlisting}

The plan materializes the daily-watch intermediate result. It scans
\texttt{ViewSegment}, searches \texttt{Impression} by integer primary
key, and uses temporary B-trees for grouping and ordering. The
self-join of the intermediate result uses a Bloom filter and an
automatically created covering index on \texttt{user\_id}.

\subsection{F12: Completed Recommendation Impressions}

F12 finds sessions in which a recommended video was shown and
watched to completion. It joins \texttt{Session}, \texttt{Turn},
\texttt{AgentRecommendation}, and \texttt{Impression}, and uses a
correlated \texttt{EXISTS} query against \texttt{ViewSegment}.

\subsubsection{Plan before index creation}

\begin{lstlisting}[style=textstyle]
|--SCAN i
|--SEARCH ar USING INTEGER PRIMARY KEY (rowid=?)
|--SEARCH t USING INTEGER PRIMARY KEY (rowid=?)
|--SEARCH s USING INTEGER PRIMARY KEY (rowid=?)
|--BLOOM FILTER ON vs (impression_id=? AND reached_end=?)
|--SEARCH vs EXISTS USING AUTOMATIC PARTIAL COVERING INDEX
|  (impression_id=? AND reached_end=?)
|--USE TEMP B-TREE FOR DISTINCT
`--USE TEMP B-TREE FOR ORDER BY
\end{lstlisting}

SQLite scans \texttt{Impression}, performs primary-key lookups in
the related tables, and uses an automatically created partial
covering index for the correlated completion check. Temporary
B-trees support \texttt{DISTINCT} and ordering.

\subsubsection{Plan after index creation}

I created the following index:

\begin{lstlisting}[style=sqlstyle]
CREATE INDEX IF NOT EXISTS
    idx_viewsegment_impression_reached
ON ViewSegment(impression_id, reached_end);
\end{lstlisting}

The resulting plan was:

\begin{lstlisting}[style=textstyle]
|--SCAN i
|--SEARCH ar USING INTEGER PRIMARY KEY (rowid=?)
|--SEARCH vs EXISTS USING COVERING INDEX
|  idx_viewsegment_impression_reached
|  (impression_id=? AND reached_end=?)
|--SEARCH t USING INTEGER PRIMARY KEY (rowid=?)
|--SEARCH s USING INTEGER PRIMARY KEY (rowid=?)
|--USE TEMP B-TREE FOR DISTINCT
`--USE TEMP B-TREE FOR ORDER BY
\end{lstlisting}

The persistent covering index is used for the correlated existence
check. The scan of \texttt{Impression} and the temporary structures
for duplicate elimination and ordering remain.

\subsection{One Full-Scan Cost Estimate}

For the required scan-cost estimate, I use the recorded full scan
of \texttt{ViewSegment} in F10.

The table contains:

\[
N = 8{,}892{,}547 \text{ rows}.
\]

For a simplified model, I assume a 4096-byte page and an average
row size of 64 bytes. Therefore, the estimated number of rows per
page is:

\[
\text{Rows per page}
=
\left\lfloor \frac{4096}{64} \right\rfloor
=64.
\]

The estimated number of pages is:

\[
B(\text{ViewSegment})
=
\left\lceil \frac{8{,}892{,}547}{64} \right\rceil
=
138{,}946.
\]

Under these assumptions, a sequential scan reads approximately
138,946 data pages. This is a simplified page-transfer estimate.
Actual SQLite I/O depends on its page layout, record sizes, indexes,
cache state, and other implementation details.

\subsection{Summary of Access-Path Findings}

\begin{table}[H]
\centering
\small
\begin{tabular}{p{1.1cm}p{5.0cm}p{5.0cm}}
\toprule
\textbf{Query} & \textbf{Before} & \textbf{After} \\
\midrule
F4 &
Scan of retraction-side \texttt{Signal}; primary-key lookup for
original signal &
Indexed search on signal type; primary-key lookup retained \\

F7 &
Scan of \texttt{Turn}; session lookup by primary key &
Date-range search on \texttt{Session}; indexed lookup into
\texttt{Turn} \\

F10 &
Full scan of \texttt{ViewSegment}; lookup into \texttt{Impression} &
Plan inspected for high-volume join and aggregation \\

F12 &
Scan of \texttt{Impression}; automatic partial covering index for
completion check &
Persistent covering index for the correlated completion check \\
\bottomrule
\end{tabular}
\caption{Summary of recorded access paths.}
\end{table}

% ============================================================
\section{Task 2: Composite-Index Experiments}
% ============================================================

\subsection{Index 1: Signal Retractions}

The first index is:

\begin{lstlisting}[style=sqlstyle]
CREATE INDEX IF NOT EXISTS
    idx_signal_type_retracts_created
ON Signal(signal_type, retracts_signal_id, created_at);
\end{lstlisting}

The leading column, \texttt{signal\_type}, supports the equality
conditions that identify likes and retractions. The next column,
\texttt{retracts\_signal\_id}, stores the reference to the signal
being retracted. The timestamp is also stored in the index.

The recorded plan changed from \texttt{SCAN r} to
\texttt{SEARCH r USING COVERING INDEX}. The plan still contains a
temporary B-tree for ordering, so the index does not eliminate all
query-processing work.

\subsection{Index 2: Session Date Range}

The second index is:

\begin{lstlisting}[style=sqlstyle]
CREATE INDEX IF NOT EXISTS
    idx_session_started_session
ON Session(started_at, session_id);
\end{lstlisting}

The leading \texttt{started\_at} column supports the date-range
predicate in F7. The \texttt{session\_id} column is included so
that SQLite can obtain the session identifier directly from the
index. The recorded plan uses a covering index for the range search.

The plan also changes the join sequence: SQLite starts with
sessions in the selected date range and then searches for their
turns. This is consistent with reducing the set of sessions before
processing subsequent joins.

\subsection{Index 3: View Completion Check}

The third index is:

\begin{lstlisting}[style=sqlstyle]
CREATE INDEX IF NOT EXISTS
    idx_viewsegment_impression_reached
ON ViewSegment(impression_id, reached_end);
\end{lstlisting}

F12 checks whether an impression has a \texttt{ViewSegment} row
with \texttt{reached\_end = 1}. Both columns in the correlated
condition are included in the index. The recorded plan uses this
persistent covering index instead of the automatically created
partial covering index that appeared in the baseline plan.

\subsection{Rejected Index}

I considered but did not build the following index:

\begin{lstlisting}[style=sqlstyle]
-- Not created:
-- CREATE INDEX idx_impression_recommendation_video_user
-- ON Impression(recommendation_id, video_id, user_id);
\end{lstlisting}

This index could support the matching conditions between
\texttt{Impression} and \texttt{AgentRecommendation} in F12.
However, \texttt{Impression} contains 15 million rows. Building
another index on this table would require additional storage and
index-build time. I left it uncreated rather than claiming a
performance improvement that I did not measure.

\subsection{Runtime Evidence and Limitations}

The before-and-after execution plans were captured for the three
index experiments. F4 changed from a scan of the retraction-side
\texttt{Signal} alias to an indexed search. F7 changed from a scan
of \texttt{Turn} to an indexed date-range search on
\texttt{Session}. F12 changed from an automatically created
partial covering index to the persistent covering index for the
completion check.

The complete queries were timed once with the target index absent
and once after recreating it. The measured wall-clock times were:

\begin{center}
\begin{tabular}{lrrr}
\toprule
Query & Before (s) & After (s) & Observed reduction (\%) \\
\midrule
F4: Signal retractions & 0.09 & 0.03 & 66.7 \\
F7: Session cost & 0.13 & 0.08 & 38.5 \\
F12: Completed recommendations & 6.94 & 3.46 & 50.1 \\
\bottomrule
\end{tabular}
\end{center}

The observed reduction is calculated as
$(T_{\mathrm{before}}-T_{\mathrm{after}})/T_{\mathrm{before}}
\times 100$. These are single-run measurements, not controlled
benchmarks. Cache state and system load may affect the results, so
the reductions are preliminary observations rather than guaranteed
performance improvements. The execution plans provide complementary
evidence of changes in access paths.

% ============================================================
\section{Task 3: Join Strategy and Cost Model}
% ============================================================

\subsection{Queries Joining High-Volume Tables}

I used F10 and F12 to examine joins involving the large
\texttt{ViewSegment} and \texttt{Impression} tables. F10 aggregates
watch information across these tables. F12 matches impressions
to recommendations and checks completion through
\texttt{ViewSegment}.

The recorded plans contain scans, indexed searches, primary-key
lookups, and temporary structures. SQLite uses nested-loop-style
join processing rather than a native hash join or sort-merge join.
In these plans, one input is scanned or searched and matching rows
are retrieved from the other input through indexes or primary-key
lookups.

This distinction matters because the cost of a nested-loop join
depends heavily on the access path for the inner input. An indexed
lookup is not equivalent to repeatedly scanning the entire inner
table.

\subsection{Cardinality and Page Estimates}

For the simplified cost comparison, I use the recorded row counts
and assume 64 rows per page:

\[
B(R) =
\left\lceil\frac{8{,}892{,}547}{64}\right\rceil
=138{,}946
\]

for \texttt{ViewSegment}, and

\[
B(S) =
\left\lceil\frac{15{,}000{,}000}{64}\right\rceil
=234{,}375
\]

for \texttt{Impression}.

These estimates use the same simplified page-size assumptions
as Task 1. They are not the actual number of SQLite pages reported
by the database.

\subsection{Block Nested-Loop Join}

Let \(M=100\) be the assumed number of buffer pages. The textbook
block nested-loop join estimate is:

\[
B(R) +
\left\lceil\frac{B(R)}{M-2}\right\rceil B(S).
\]

Substituting the estimates gives:

\[
138{,}946 +
\left\lceil\frac{138{,}946}{98}\right\rceil
(234{,}375).
\]

Since

\[
\left\lceil\frac{138{,}946}{98}\right\rceil=1{,}418,
\]

the estimated number of page transfers is:

\[
138{,}946 + 1{,}418(234{,}375)
=332{,}482{,}696.
\]

This is the cost of the simplified block nested-loop model under
the stated assumptions. It is not a measured cost of the actual
SQLite execution.

\subsection{Sort-Merge Join}

Under the simplified two-pass model, a sort-merge join can be
estimated at approximately:

\[
3(B(R)+B(S)).
\]

Therefore:

\[
3(138{,}946+234{,}375)
=1{,}119{,}963
\]

page transfers.

This estimate assumes that the inputs can be sorted in the
available memory and that the relevant two-pass assumptions hold.
If the data do not fit those assumptions, additional passes may
be required.

\subsection{Hash Join}

Under the same simplified two-pass assumptions, a Grace hash join
has an approximate cost of:

\[
3(B(R)+B(S))
=1{,}119{,}963
\]

page transfers.

This estimate assumes that the partitions can be processed within
the available memory and that the hash-join algorithm's assumptions
hold. It is a theoretical comparison, not a claim that SQLite used
a hash join.

\begin{table}[H]
\centering
\begin{tabular}{lr}
\toprule
\textbf{Join model} & \textbf{Estimated page transfers} \\
\midrule
Block nested-loop & 332,482,696 \\
Two-pass sort-merge & 1,119,963 \\
Two-pass Grace hash & 1,119,963 \\
\bottomrule
\end{tabular}
\caption{Theoretical cost comparison under the stated assumptions.}
\end{table}

Under this simplified model, sort-merge and hash join have much
lower estimated transfer costs than block nested-loop join for
these input sizes. However, SQLite's actual plans use the join
strategies available in SQLite and the indexes present in the
database. The textbook estimates do not prove that the actual
SQLite plan is more expensive than an unavailable algorithm would
be in practice.


```latex
\subsection{Three-Table Join Ordering}

Consider a query joining \texttt{Session}, \texttt{Turn}, and
\texttt{PromptTemplateVersion}, with a date-range predicate on
\texttt{Session.started\_at}. A reasonable join order is to filter
\texttt{Session} first, join the qualifying sessions to \texttt{Turn},
and then retrieve the corresponding template versions.

The measured cardinalities support this choice. Filtering
\texttt{Session} to August 2026 produces 15,886 rows. Joining these
sessions to \texttt{Turn} produces 43,779 rows. By comparison, joining
the complete \texttt{Turn} relation to \texttt{PromptTemplateVersion}
without applying the date filter produces 276,011 rows. The complete
three-table query with the August predicate returns 43,779 rows.

The SQLite execution plan for the date-filtered query is:

\begin{samepage}
\small
\begin{lstlisting}
SEARCH s USING COVERING INDEX idx_session_started_session
    (started_at>? AND started_at<?)
SEARCH t USING INDEX sqlite_autoindex_Turn_1 (session_id=?)
SEARCH ptv USING INTEGER PRIMARY KEY (rowid=?)
\end{lstlisting}
\end{samepage}

The plan uses the date-range index to find qualifying sessions,
searches \texttt{Turn} using the session ID, and retrieves matching
template versions using their integer primary key. This is consistent
with the proposed filtered-session-first strategy.

These measurements support the proposed join order but do not prove
that it is globally cheapest. The 276,011-row measurement represents
the unfiltered \texttt{Turn}-to-\texttt{PromptTemplateVersion} join,
whereas the 43,779-row result includes the August session predicate.
A definitive comparison of all candidate orders would require
comparable intermediate-cardinality estimates and cost measurements
for each alternative.

% ============================================================
\section{Task 4: Predicate Pushdown}
% ============================================================

\subsection{Query and Selective Predicate}

I used F7, which calculates session-level inference costs for
sessions that started during August 2026. Its selective predicate
restricts the session relation to the desired date interval:

\begin{lstlisting}[style=sqlstyle]
s.started_at >= '2026-08-01T00:00:00Z'
AND s.started_at < '2026-09-01T00:00:00Z'
\end{lstlisting}

The predicate refers only to the \texttt{Session} relation. This
makes it suitable for an explicit predicate-pushdown rewrite
under the query's inner-join semantics.

\subsection{Naive Relational-Algebra Plan}

Let \(S\), \(T\), \(V\), and \(P\) denote \texttt{Session},
\texttt{Turn}, \texttt{PromptTemplateVersion}, and
\texttt{ModelPricingPeriod}, respectively. Let \(\sigma_d\)
denote the selection for the August date range.

A naive logical plan can be represented as:

\[
\sigma_d(S \bowtie T \bowtie V \bowtie P).
\]

This places the selection after the joins in the logical expression.
If implemented literally, it can produce intermediate join results
containing sessions that are subsequently discarded by the date
predicate.

\subsection{Pushed-Down Relational-Algebra Plan}

The pushed-down expression is:

\[
(\sigma_d(S)) \bowtie T \bowtie V \bowtie P.
\]

This applies the date predicate to \texttt{Session} before the
joins. The intended benefit is to reduce the number of session
tuples passed into the subsequent join operations.

\subsection{Equivalence Rule and Limitation}

For inner joins, a selection predicate that refers only to
attributes of one input relation can be pushed below the join
without changing the result, provided the join conditions and
predicate semantics remain unchanged.

This rule must be applied carefully to outer joins. Moving a
predicate across an outer join can change the result, particularly
when the predicate filters out NULL-extended rows that the outer
join was meant to preserve.

\subsection{Original SQL}

\begin{lstlisting}[style=sqlstyle]
SELECT
    s.session_id,
    ptv.template_version_id,
    ptv.version_number,
    ROUND(
        SUM(
            (t.input_tokens - t.cached_input_tokens) * p.input_rate
            + t.cached_input_tokens * p.cached_input_rate
            + t.output_tokens * p.output_rate
        ),
        6
    ) AS session_cost
FROM Session s
JOIN Turn t
    ON t.session_id = s.session_id
JOIN PromptTemplateVersion ptv
    ON ptv.template_version_id =
       t.prompt_template_version_id
JOIN ModelPricingPeriod p
    ON p.model_id = t.model_id
   AND t.created_at >= p.valid_from
   AND (p.valid_to IS NULL OR t.created_at < p.valid_to)
WHERE s.started_at >= '2026-08-01T00:00:00Z'
  AND s.started_at < '2026-09-01T00:00:00Z'
GROUP BY
    s.session_id,
    ptv.template_version_id,
    ptv.version_number
HAVING session_cost > 0.05
ORDER BY
    session_cost DESC,
    s.session_id,
    ptv.version_number;
\end{lstlisting}

\subsection{Rewritten SQL}

I made the early filtering explicit using a common table expression
named \texttt{filtered\_sessions}:

\begin{lstlisting}[style=sqlstyle]
WITH filtered_sessions AS (
    SELECT session_id, user_id
    FROM Session
    WHERE started_at >= '2026-08-01T00:00:00Z'
      AND started_at < '2026-09-01T00:00:00Z'
)
SELECT
    s.session_id,
    ptv.template_version_id,
    ptv.version_number,
    ROUND(
        SUM(
            (t.input_tokens - t.cached_input_tokens) * p.input_rate
            + t.cached_input_tokens * p.cached_input_rate
            + t.output_tokens * p.output_rate
        ),
        6
    ) AS session_cost
FROM filtered_sessions s
JOIN Turn t
    ON t.session_id = s.session_id
JOIN PromptTemplateVersion ptv
    ON ptv.template_version_id =
       t.prompt_template_version_id
JOIN ModelPricingPeriod p
    ON p.model_id = t.model_id
   AND t.created_at >= p.valid_from
   AND (p.valid_to IS NULL OR t.created_at < p.valid_to)
GROUP BY
    s.session_id,
    ptv.template_version_id,
    ptv.version_number
HAVING session_cost > 0.05
ORDER BY
    session_cost DESC,
    s.session_id,
    ptv.version_number;
\end{lstlisting}

Both versions retain the same date range, join conditions, cost
calculation, grouping, \texttt{HAVING} condition, and ordering.
The rewritten query makes the early selection explicit.

\subsection{Execution-Plan Evidence}

Both the original and rewritten queries produced the same
execution-plan operations. SQLite uses the date-range index
\texttt{idx\_session\_started\_session} to search qualifying
sessions, looks up turns by session ID, retrieves template
versions using their integer primary key, and searches pricing
periods using the relevant index. Both plans also use temporary
B-trees for grouping and ordering.

The recorded plan for the original query is:

\begin{lstlisting}[style=textstyle]
SEARCH s USING COVERING INDEX idx_session_started_session
    (started_at>? AND started_at<?)
SEARCH t USING INDEX sqlite_autoindex_Turn_1 (session_id=?)
SEARCH ptv USING INTEGER PRIMARY KEY (rowid=?)
SEARCH p USING INDEX sqlite_autoindex_ModelPricingPeriod_1
    (model_id=? AND valid_from<?)
USE TEMP B-TREE FOR GROUP BY
USE TEMP B-TREE FOR ORDER BY
\end{lstlisting}

\subsection{Runtime Comparison}

I executed both queries five times, alternating their execution
order. Each execution returned 38,005 rows, and the ordered
results were identical in every round.

The median execution times were:

\begin{center}
\begin{tabular}{lr}
\hline
Query & Median runtime (seconds) \\
\hline
Original SQL & 0.0790 \\
Rewritten SQL & 0.0784 \\
\hline
\end{tabular}
\end{center}

The rewritten query was approximately 0.76\% faster based on
these median measurements. This difference is too small to
establish a meaningful performance improvement from the rewrite.
The first execution of the original query took 0.2910 seconds,
whereas subsequent executions were much faster, indicating that
warm-up or caching effects may influence the measurements.

The execution plans and timing results suggest that SQLite's
optimizer already applies the beneficial filtering strategy to
the original query. Therefore, making the predicate pushdown
explicit using a common table expression does not produce a
meaningful observed improvement in this experiment.

% ============================================================
\section{Task 5: Concurrency, Schedules, and Isolation}
% ============================================================

\subsection{A1 Deliverable G.3 (T2): Read Isolation}

In the A1 T2 demonstration, Connection 1 began a transaction and
inserted an uncommitted moderation decision for video 1 with state
\texttt{live}. Connection 1 then queried the current state of the
video and could observe its own uncommitted change. Before
Connection 1 committed, Connection 2 queried the same video and
observed the previously committed \texttt{pending} state.
Connection 1 subsequently rolled back.

The relevant operation sequence is:

\begin{table}[H]
\centering
\small
\begin{tabular}{clp{8cm}}
\toprule
\textbf{Order} & \textbf{Connection} & \textbf{Operation} \\
\midrule
1 & C1 & Begin transaction. \\
2 & C1 & Insert an uncommitted \texttt{live} moderation decision
for video 1. \\
3 & C1 & Read the current state and observe \texttt{live}. \\
4 & C2 & Read the current state and observe the committed
\texttt{pending} state. \\
5 & C1 & Roll back the transaction. \\
\bottomrule
\end{tabular}
\caption{Summary of the A1 T2 read-isolation experiment.}
\end{table}

This experiment demonstrates that Connection 2 does not see
Connection 1's uncommitted insert. Because Connection 1 rolls back,
the experiment does not establish a precedence edge between two
committed transactions that would demonstrate a non-serializable
schedule.

The corresponding precedence graph for the committed operations
has no conflict edge between these two transactions. It is
therefore acyclic for this demonstrated interleaving. This is a
read-isolation experiment rather than evidence of a committed
write-skew anomaly.

\subsection{Proposed Write-Skew Scenario}

I considered an application invariant involving videos 1 and 2:
at least one of these videos must remain in the \texttt{live} state.

The proposed scenario has two transactions. Each reads the latest
moderation state of both videos. Transaction \(T_1\) then decides
to update video 1, while transaction \(T_2\) decides to update
video 2. If both transactions base their decisions on the same
original state, each may assume that the other video will remain
live.

The abstract schedule is:

\begin{table}[H]
\centering
\small
\begin{tabular}{clp{7.8cm}}
\toprule
\textbf{Order} & \textbf{Transaction} & \textbf{Operation} \\
\midrule
1 & \(T_1\) & Read state of video 1. \\
2 & \(T_2\) & Read state of video 1. \\
3 & \(T_1\) & Read state of video 2. \\
4 & \(T_2\) & Read state of video 2. \\
5 & \(T_1\) & Write a new moderation decision for video 1. \\
6 & \(T_2\) & Write a new moderation decision for video 2. \\
\bottomrule
\end{tabular}
\caption{Abstract schedule for the proposed write-skew scenario.}
\end{table}

The conflicting operations yield two edges in the precedence graph:

\begin{itemize}[leftmargin=2em]
    \item \(T_2 \rightarrow T_1\): \(T_2\)'s read of video 1
    occurs before \(T_1\)'s write to video 1.
    \item \(T_1 \rightarrow T_2\): \(T_1\)'s read of video 2
    occurs before \(T_2\)'s write to video 2.
\end{itemize}

The graph contains a cycle:

\[
T_1 \rightarrow T_2 \rightarrow T_1.
\]

Therefore, the abstract schedule is not conflict-serializable.
If both transactions committed their independent updates, both
videos could be demoted, violating the assumed invariant. This
would be a write-skew anomaly.

However, the schedule above is an abstract concurrency analysis.
It is not evidence that SQLite committed both writes. SQLite
allows only one writer at a time, and in WAL mode a transaction
that tries to write from an outdated read snapshot may fail with
\texttt{SQLITE\_BUSY\_SNAPSHOT}. I did not record a successful
committed execution of this anomaly.

\subsection{Two-Connection Test of \texttt{BEGIN IMMEDIATE}}

I tested SQLite's writer reservation using two independent Python
SQLite connections against a disposable copy of the scaled database,
named \texttt{concurrency\_test.db}. Both connections used a
one-second timeout.

Connection 1 executed \texttt{BEGIN IMMEDIATE} and acquired the
writer reservation. While Connection 1 held it, Connection 2
attempted \texttt{BEGIN IMMEDIATE}. The observed output was:

\begin{lstlisting}[style=textstyle]
Connection 1: BEGIN IMMEDIATE
Connection 1: writer reservation acquired
Connection 2: attempting BEGIN IMMEDIATE
Connection 2 result: OperationalError: database is locked
Connection 1: ROLLBACK complete
\end{lstlisting}

Connection 1 rolled back successfully. This confirms that the
second connection could not acquire the writer reservation during
this test. 

```latex
\subsection{Observed Lost Update and Prevention}

To demonstrate an actual concurrency anomaly, I created a small,
separate SQLite database named \texttt{concurrency\_demo\_small.db}.
The table contained one account with an initial balance of 100.
Two independent connections each attempted to deduct 10 from the
balance.

In the unsafe experiment, both connections read the balance before
either updated it. Each calculated the new balance from its own
read value and wrote 90. Consequently, the second write overwrote
the first deduction.

The observed output was:

\begin{lstlisting}[style=textstyle]
PART A: UNSAFE READ-MODIFY-WRITE
T1 reads balance=100
T2 reads balance=100
T2 writes balance=90
T1 writes balance=90
Expected balance if both deductions apply: 80
Actual final balance: 90
\end{lstlisting}

The final balance of 90 demonstrates a lost update: both
transactions completed their deductions logically, but only one
deduction was reflected in the stored balance.

I then repeated the experiment with \texttt{BEGIN IMMEDIATE}
issued before each transaction read the balance. The transaction
that acquired the writer reservation first read and updated the
balance and committed. The second transaction then read the
updated balance and applied its deduction.

The observed output was:

\begin{lstlisting}[style=textstyle]
PART B: BEGIN IMMEDIATE
T1 reads balance=100
T1 commits balance=90
T2 reads balance=90
T2 commits balance=80
Expected final balance: 80
Actual final balance: 80

PASS: Both deductions applied under transaction locking.
\end{lstlisting}

The final balance of 80 confirms that both deductions were
preserved in the protected experiment. This demonstrates the
prevention of the observed lost update when the complete
read-modify-write operation is performed inside a transaction
that acquires its writer reservation before reading.

This experiment uses a separate account table to isolate the
concurrency behaviour. It demonstrates a lost update and its
prevention, rather than a committed write-skew anomaly in the
ScrollSense moderation schema.
```

\subsection{Prevention and Trade-Off}

The experiment demonstrates that \texttt{BEGIN IMMEDIATE} can
prevent the observed lost update when the transaction acquires
the writer reservation before reading and retains it through the
update and commit. The earlier two-connection locking test also
showed that a second connection could not acquire the writer
reservation while the first connection held it.

The trade-off is that competing writers must serialize, which
can increase waiting time and reduce write concurrency. If a
connection cannot acquire the writer reservation within its
configured timeout, SQLite may return \texttt{SQLITE\_BUSY};
the application may then need to retry the transaction. The
transaction should be retried as a whole so that its reads and
decisions are made against current data.

% ============================================================
\section{Conclusion}
% ============================================================

In this assignment, I examined the query plans of several ScrollSense
queries on a scaled SQLite database. The index experiments changed
the access paths used for F4, F7, and F12. In particular, SQLite
used the signal-type index for the retraction query, the session
date-range index for F7, and the persistent covering index for
F12's correlated completion check.

I also compared textbook join-cost estimates. Under the simplified
page model, block nested-loop join has a much larger estimated
transfer cost than two-pass sort-merge or Grace hash join for the
selected input sizes. These estimates are theoretical and do not
represent measured SQLite costs, because SQLite's actual execution
uses the join strategies available to it and the indexes in the
database.

For predicate pushdown, I rewrote F7 so that the date filter is
explicitly applied to \texttt{Session} before the joins. The captured
execution plan already shows an indexed range search on
\texttt{Session}, so the plan supports early filtering but does
not establish an additional speedup from the rewrite.

Finally, the concurrency experiments demonstrated a lost update
under an unsafe read-modify-write sequence: the final balance
was 90 instead of the expected 80. Repeating the experiment with
\texttt{BEGIN IMMEDIATE} preserved both deductions and produced
the expected final balance of 80. The separate write-skew schedule
has a cyclic precedence graph, but a committed write-skew outcome
was not observed in SQLite. These results distinguish the observed
lost-update anomaly from the theoretical write-skew analysis.

The three index experiments have recorded before-and-after
wall-clock measurements: F4 took 0.09 seconds before and 0.03
seconds after indexing; F7 took 0.13 and 0.08 seconds; and F12
took 6.94 and 3.46 seconds. These are single-run observations,
not controlled benchmarks, and may be affected by cache state
and system load. A directly comparable timing experiment for the
original F7 SQL and the explicit predicate-pushdown rewrite
remains outstanding.

% ============================================================
\section*{Appendix: Reproducibility and Assistance}
% ============================================================

The working directory contains the SQL queries, schema, view
definitions, index definitions, data-generation script, transaction
demonstration, and analysis notes used in this work. The scaled
database was kept separate from the original Assignment 1 files.
The concurrency test was performed on a disposable database copy.

I used an AI assistant for help with understanding execution plans,
structuring explanations, and editing the report. The final report
should be read alongside the SQL files and the recorded experimental
evidence. Any code, interpretation, and results submitted remain
my responsibility to verify against the assignment requirements.

\end{document}
